\BourbakiExerciseGroup

\begin{enumerate}
\def\labelenumi{\arabic{enumi})}
\item
  设 A 为主理想整环，设 M 为挠 A-模，设 x, y 为 M 的两个元素，并设 Aα 与 Aβ 分别为 x, y 的零化子；若 δ 为 α 与 β 的一个最大公因子，证明 \(x + y\) 的零化子 Aγ 使得 γ 整除 α 与 β 的最小公倍数 αβ/δ，并且是 αβ/δ² 的倍数。当 α 与 β 是 A 的同一个不可约元素 π 的不同幂 \(\pi^{\mu}\) , \(\pi^{\nu}\) 时，证明 \(\gamma = \pi^{\sup(\mu, \nu)}\) 。
\item
  设 A 为主理想整环，设 \(\pi\) 为 A 的一个不可约元素，并设 M 为 \(\pi\) -准素模；证明对所有 \(x \in M\) 以及所有 \(\alpha \in A\) （不是 \(\pi\) 的倍数），存在唯一的 \(y \in M\) 使得 \(x = \alpha y\) （利用贝祖恒等式）；令 \(y = \alpha^{-1}x\) 。由此推出，如果 \(A_{(\pi)}\) 表示由 A 的分式域中形如 \(\beta/\alpha\) 、满足 \(\alpha, \beta \in A\) 且 \(\alpha\) 不是 \(\pi\) 的倍数的元素构成的主理想整环（VII，第48页，习题4），则 M 上存在唯一的 \(A_{(\pi)}\) -模结构，使得算子环到 A 的限制给出 M 上原有的 A-模结构。这样定义的 \(A_{(\pi)}\) -模被称为与 M 典范关联；其子模与 M 的 A-子模完全相同。
\item
  设 A 为主理想整环。
\end{enumerate}

\begin{enumerate}
\def\labelenumi{\alph{enumi})}
\item
  一个 A-模 M 是内射的（II，第389页，习题11），当且仅当对于所有 \(x \in \mathbf{M}\) 以及 A 中所有 \(\alpha \neq 0\) ，存在 \(y \in \mathbf{M}\) 使得 \(x = \alpha y\) ；此时称 A-模 M 是可除的。若 E 是任意 A-模，则 E 的所有可除子模之和是 E 的最大可除子模 \(\Delta(E)\) 。对 E 的每个子模 F，我们有 \(\Delta(F) \subset \Delta(E)\) 。
\item
  设 M 是一个可除 A-模，设 \(\pi\) 为 A 的一个不可约元素，且令 \(x_{0}\) 是 M 中的一个元素，其零化子具有形式 \(A\pi^{n}\) (n \textgreater{} 0 为整数）。递归地定义 M 的元素的一个序列 \((x_{s})\) ，条件为 \(x_{s} = \pi x_{s+1}\) 对每个整数 \(s \geqslant 0\) 成立。证明由序列 \((x_{s})\) 生成的 M 的子模同构于模 U = K/A 的 \(\pi\) -准素分量 \(U_{\pi}\) ，其中 K 是 A 的分式域。c) 证明 \(U_{\pi}\) 是一个可除模，并且 \(U_{\pi}\) 的每个不同于 \(U_{\pi}\) 和 0 的子模都是循环的，且由 \(\pi\) 的某个幂关于 mod A 的类生成；由此推出 \(U_{\pi}\) 是不可分解的（II，第393页，习题21）。
\end{enumerate}

\(d\) ）证明每个不可分解的可除 \(\mathbf{A}\) -模同构于 \(\mathbf{K}\) ，或同构于诸模 \(\mathrm{U}_{\pi}\) 之一（利用 \(a\) ）和 \(b)\) ）。

\begin{enumerate}
\def\labelenumi{\alph{enumi})}
\setcounter{enumi}{4}
\tightlist
\item
  证明每个可除A-模都同构于一族不可分解可除子模的直和（将Zorn引理应用于其和为直和的不可分解可除子模族）。
\end{enumerate}

\begin{enumerate}
\def\labelenumi{\arabic{enumi})}
\setcounter{enumi}{3}
\tightlist
\item
  一个族 \((x_{i})\) （由模 \(\mathbf{M}\) ——主理想整环 \(\mathbf{A}\) 上的——的非零元素组成）称为拟自由的，如果每个形如 \(\sum_{i} \alpha_{i} x_{i} = \beta y\) 的关系式 ( \(y \in \mathbf{M}, \alpha_{i} \in \mathbf{A}\) ,
\end{enumerate}

\(\beta \in \mathbf{A}\) ) 都蕴含存在一族元素 \(\alpha_{i}^{\prime} \in \mathbf{A}\) ，使得 \(\alpha_{i}x_{i} = \beta \alpha_{i}^{\prime} x_{i}\) 对每个指标 \(i\) 成立。于是关系式 \(i \neq k\) 蕴含 \(x_{i} \neq x_{k}\) ；拟自由族的元素组成的集合称为 \(\mathbf{M}\) 的拟自由子集。

\begin{enumerate}
\def\labelenumi{\alph{enumi})}
\item
  设 \(\pi\) 是 A 的一个不可约元素。证明：在具有零化子 \(\mathbf{A}\pi^n\) ( \(n \geqslant 1\) ) 的 A-模中，任何具有零化子 \(\mathbf{A}\pi^n\) 的元素都是拟自由的（使用贝祖恒等式）。
\item
  设 \((x_{i})\) 是 A-模 M 中的一个拟自由族，并设 N 为由该族生成的 M 的子模。证明对 M/N 的每个元素 \(\dot{x}\) ，存在一个元素 \(x \in \dot{x}\) ，其零化子等于 \(\dot{x}\) 在 M/N 中的零化子。若 \(\dot{x}\) 在 M/N 中还是拟自由的，证明由 \((x_{i})\) 以及 x 组成的族在 M 中是拟自由的。
\item
  由 a) 和 b) 可推出，若 A-模 M 的零化子非零，则 M 是循环模的直和（化归到 \(\pi\) -准素模的情形，并对 M 的拟自由元素族应用 Zorn 引理）（参见 VII，第 19 页，定理 2）。
\end{enumerate}

\begin{enumerate}
\def\labelenumi{\arabic{enumi})}
\setcounter{enumi}{4}
\item
  设 M 是主理想整环 A 上的有限生成挠模。证明 M 是有限长度的模；因此，M 的任意非空子模集合都有极大元和极小元。
\item
  设 M 是主理想整环 A 上的挠模，它不是有限生成的，但其所有真子模都是有限生成的。证明存在 A 的一个不可约元素 \(\pi\) ，使得 M 同构于模 U = K/A 的 \(\pi\) -准素分量 \(U_{\pi}\) （习题 3，c））。（首先证明 M 的非零 \(\pi\) -准素分量不能多于一个，因而 M 是一个 \(\pi\) -准素模，对某个不可约元素 \(\pi \in A\) 。然后证明 M = \(\pi M\) ，这可通过观察 M/ \(\pi M\) 若非零则不可能有限生成得到；但 M/ \(\pi M\) 是域 A/ \(\pi A\) 上的向量空间，并且会包含一个非有限生成的真子空间。）
\item
  设 M 是主理想整环 A 上的模；M 的子模 N 被称为纯的，如果对所有 \(\alpha \in A\) ，则对 \(N \cap (\alpha M) = \alpha N\) .
\end{enumerate}

\begin{enumerate}
\def\labelenumi{\alph{enumi})}
\item
  M 的子模 N 是纯的，当且仅当对于所有 \(\alpha \in A\) 和每个元素 \(\dot{x} \in M/N\) （其零化子为 \(A\alpha\) ），存在一个元素 \(x \in \dot{x}\) ，其零化子为 \(A\alpha\) 。
\item
  证明 M 的挠子模是纯的。 M 的元素的族 \((x_{t})\) 是拟自由的（习题 4），当且仅当子模 \(Ax_{t}\) 的和是直和并且是 M 的纯子模。证明如果 P 和 Q 是 M 的子模，使得 \(P \cap Q\) 和 \(P + Q\) 是纯的，则 P 和 Q 是纯的。如果 Q 是 M 的纯子模，且 \(P \supset Q\) 是 M 的子模，则 P 在 M 中是纯的，当且仅当 P/Q 在 M/Q 中是纯的。
\item
  证明 M 的每个作为 M 的直因子的子模 N 都是纯的。对于 M 的直因子的滤过族的并集，同样成立。反之，如果 N 是 M 的纯子模，且 M/N 是循环子模的直和，证明 N 是 M 的直因子（使用 a））。给出 M 的直因子 P 和 Q 的例子，使得 \(P \cap Q\) （相应地 \(P + Q\) ）不是纯子模（对于第一个例子，取 \(M = (\mathbf{Z}/4\mathbf{Z}) \oplus (\mathbf{Z}/2\mathbf{Z})\) ，对于第二个例子，取 \(M = Z^{2}\) ）。
\item
  设 N 是 M 的纯子模，其零化子 Aα 非零。如果 f 是从 M 到 M/αM 的自然同态，证明 f 限制为从 N 到 f(N) 的同构，并且 f(N) 是 M/αM 的纯子模（使用 Bézout 恒等式）。由此推出（借助 c）和习题 4，c）） \(f(N)\) 是 M/ \(\alpha M\) 中的直因子，并且 N 是 M 中的直因子。特别地，如果模 M 的挠子模具有非零零化子（当它是有限生成时就是这种情况），则它是 M 中的直因子（参见 VII，第 20 页，定理 2 的推论 1）。
\item
  设 \(p\) 是一个素数；设 \(\mathbf{N}\) 表示循环 \(\mathbf{Z}\) -模 \(\mathbf{Z} / (p^n)\) ( \(n \geqslant 1\) ) 的直和，并令 \(e_n\) 表示 \(\mathbf{N}\) 中这样的元素：其所有坐标均为零，除了第 \(n\) 个坐标，它等于 1 模 \(p^n\) 的类。设 \(\mathbf{M}\) 为 \(\mathbf{Z}\) -模 \(\mathbf{N} \times \mathbf{Q}\) 的由元素 \((e_n, p^{-n})\) 生成的子模。证明挠子模 \(\mathbf{P}\) （ \(\mathbf{M}\) 的）不是 \(\mathbf{M}\) 中的直因子（注意， \(\mathbf{M} / \mathbf{P}\) 的每个元素都能被 \(p^n\) 整除，对所有 \(n\) ，并且 \(\mathbf{M}\) 中没有元素具有此性质，除了 \((0, 1)\) 的倍数）。
\end{enumerate}

\begin{enumerate}
\def\labelenumi{\arabic{enumi})}
\setcounter{enumi}{7}
\tightlist
\item
  设 A 为主理想整环，设 \(\pi\) 为 A 的一个不可约元素，并设 M 为 \(\pi\) -准素模。则 M 的元素 \(x\) 称为具有高度 \(n\) ，如果 \(x \in \pi^n\mathbf{M}\) 且 \(x \notin \pi^{n+1}\mathbf{M}\) 。若 \(x \in \pi^n\mathbf{M}\) 对所有 \(n \geqslant 1\) 成立，则称 \(x\) 具有无限高度。
\end{enumerate}

\begin{enumerate}
\def\labelenumi{\alph{enumi})}
\item
  证明M是可除的当且仅当其所有元素都具有无限高度（习题3）。
\item
  设 x 为 M 中高度为 n 的元素，使得 \(\pi x = 0\) ；证明若 \(x = \pi^{n}y\) 则 M 的循环子模 Ay 是 M 的一个直和因子（注意 Ay 是纯子模，并应用习题 7，d））。
\item
  证明一个不可除的 \(\pi\) -准素模 M 总是包含至少一个非零循环直和因子（若 \(x \in \mathbf{M}\) 具有有限高度，且存在整数 \(m\) 使得 \(\pi^m x \neq 0\) 具有无限高度，则令 \(s\) 为满足该性质的最小整数；写出 \(\pi^s x = \pi y\) ，其中 \(y\) 的高度严格大于 \(\pi^{s-1}x\) ，并将 \(b\) ) 应用于 \(y - \pi^{s-1}x\) ）。由此推出，任何不可分解的 \(\pi\) -准素模要么是循环的，要么同构于一个模 \(\mathrm{U}_{\pi}\) （习题3）。d) 证明：若 \(\pi\) -准素模 M 的每个纯子模都是 M 中的直和因子，则 M 是一个可除模 P 与一个具有非零零化子的 \(\pi\) -准素模 Q 的直和。（化归到 M 不含非零可除子模的情形，并用反证法论证。利用 b) 证明：将存在 M 中高度为 1 的元素组成的无穷序列 \((x_k)\) ，使得：1) \(x_k\) 的零化子是 \(\mathrm{A}\pi^{n_k}\) ，其中整数序列 \((n_k)\) 是严格递增的；并且 2）诸 \(\mathrm{A}x_k\) 的和是直和，且这些子模中任意有限个的和是 M 的一个直和因子。然后证明 M 的由元素 \(x_k - \pi^{n_{k+1} - n_k}x_{k+1}\) 生成的子模 N 是纯的，但不是 M 的直因子，论证如习题 7 e)）。
\end{enumerate}

\begin{enumerate}
\def\labelenumi{\arabic{enumi})}
\setcounter{enumi}{8}
\tightlist
\item
  \begin{enumerate}
  \def\labelenumii{\alph{enumii})}
  \tightlist
  \item
    设 M 为一个 \(\pi\) -准素模，并设 \(M_{0}\) 为 M 中由无限高度元素组成的子模。证明 \(\pi\) -准素模 \(M/M_{0}\) 不包含任何具有无限高度的非零元素。
  \end{enumerate}
\end{enumerate}

\begin{enumerate}
\def\labelenumi{\alph{enumi})}
\setcounter{enumi}{1}
\tightlist
\item
  设 \(p\) 是一个素数，且令 \(E\) 是 \(p\) -群 \(E_n = \mathbf{Z} / (p^n)\) ( \(n \geqslant 1\) ) 的直和；令 \(e_n\) 表示 \(1 \mod p^n\) 在 \(E_n\) 中的剩余类。设 \(H\) 是 \(E\) 的由元素 \(e_1 - p^{n-1}e_n\) （对所有整数 \(n \geqslant 2\) ）生成的子模。设 \(M\) 是商模 \(E / H\) 。证明由无限高度元素构成的子模 \(M_0\) （ \(M\) 中的）是子模 \((E_1 + H) / H\) ，它同构于 \(E_1\) ，并证明 \(M_0\) 中不存在在 \(M_0\) 中具有无限高度的非零元素。
\end{enumerate}

¶ 10) a) 设 M 是一个 π-准素模，且没有非零的无限高度元素，并设 N 是 M 的纯子模，由 M 的一个极大拟自由元素族生成（习题 4 和 7，b)）。证明 M/N 是一个可除模（用反证法论证，依次应用习题 8，c)、习题 7，c)、习题 7，b)，最后应用习题 4，c））。

\begin{enumerate}
\def\labelenumi{\alph{enumi})}
\setcounter{enumi}{1}
\item
  考虑 M 上如下定义的可度量化拓扑 \(\mathcal{T}\) ：取子模 \(\pi^n\mathbf{M}\) （对所有整数 \(n\geqslant 0\) ）作为 0 的邻域基（一般拓扑，III，第5页）。那么，M 关于子模 P 的商 M/P 是可除的，当且仅当 P 在拓扑 \(\mathcal{T}\) 下在 M 中处处稠密。
\item
  设 \(P_{1}\) , \(P_{2}\) 是 M 的两个子模，它们是纯的，在拓扑 T 下在 M 中处处稠密，并且是循环子模的直和。证明 \(P_{1}\) 和 \(P_{2}\) 是同构的（注意 \(P_{1}/\pi^{n}P_{1}\) 对所有 n 同构于 \(M/\pi^{n}M\) ）。
\item
  设 \(\hat{\mathbf{M}}\) 是 \(\mathbf{M}\) 关于拓扑 \(\mathcal{T}\) 的完备化（一般拓扑学，III，第24页）。证明 \(x \mapsto \alpha x\) 是从 \(\hat{\mathbf{M}}\) 到 \(\alpha \hat{\mathbf{M}}\) 上的严格态射，对所有 \(\alpha \in \mathbf{A}\) ， \(\pi^n \hat{\mathbf{M}}\) 是开且闭的子模，且是 \(\pi^n \mathbf{M}\) 的闭包（对所有 \(n\) ），并且商模 \(\hat{\mathbf{M}} / \pi^n \hat{\mathbf{M}}\) 与 \(\mathbf{M} / \pi^n \mathbf{M}\) 是同构的。
\item
  设 \(\bar{\mathbf{M}}\) 为 \(\hat{\mathbf{M}}\) 的挠子模；证明 \(\bar{\mathbf{M}}\) 是一个 \(\pi\) -准素模，有 \(\bar{\bar{\mathbf{M}}} = \bar{\mathbf{M}}\) ，并且 \(\mathbf{M}\) 是 \(\bar{\mathbf{M}}\) 的纯子模。
\item
  假设 \(\pi\) -准素模 M 是循环子模的一个族 \((E_{\alpha})\) 的直和。那么 T 是离散的当且仅当 M 的零化子是 0。在这种情况下，令 E 为积模 \(\prod_{\alpha} E_{\alpha}\) ，并令 F 为 E 的由这样的元素组成的子模：这些元素除
\end{enumerate}

可数多个坐标外所有坐标均为零；证明 \(\hat{\mathbf{M}}\) 随后可以（作为非拓扑模）与 \(F\) 的由这样的 \((x_{\alpha})\) 组成的子模等同：对于每个整数 \(n\) ，只有有限多个指标 \(\alpha\) 使得 \(x_{\alpha} \notin \pi^{n} E_{\alpha}\) 。然后证明这四个模 \(\mathbf{M}\) , \(\hat{\mathbf{M}}\) , \(\bar{\mathbf{M}}\) 和 \(F\) 是不同的。由此推出 \(\bar{\mathbf{M}}\) 不是循环模的直和（使用 \(e\) ），并且它不是不可分解模的直和（注意 \(\bar{\mathbf{M}}\) 不含无限高度的元素，并应用习题8，c））。

\begin{enumerate}
\def\labelenumi{\arabic{enumi})}
\setcounter{enumi}{10}
\tightlist
\item
  设 M 为一个 \(\pi\) -准素模，设N是M的一个子模，其非零元素在M中的高均为 \(\leqslant h\) 。
\end{enumerate}

\begin{enumerate}
\def\labelenumi{\alph{enumi})}
\item
  设 P 是 M 的一个纯子模，使得 \(P(\pi) \subset N(\pi)\) 且 \(P(\pi) \neq N(\pi)\) 。设 a 是 \(N(\pi)\) 中不属于 \(P(\pi)\) 的一个元素，并设 \(x_{0}\) 为 \(a + P(\pi)\) 中的一个元素，其在 M 中的高度 k 尽可能大；设 \(x_{0} = \pi^{k} y_{0}\) 。证明和 \(P + Ay_{0}\) 是直和，并且是 M 的一个纯子模。
\item
  设 B 是 M 的一个拟自由子集，使得若 P 是由 B 生成的纯子模，则 \(\mathrm{P}(\pi) \subset \mathrm{N}(\pi)\) 。证明存在 M 的一个拟自由子集 \(C \supset B\) ，使得若 Q 是由 C 生成的 M 的纯子模，则 \(\mathrm{Q}(\pi) = \mathrm{N}(\pi)\) （应用 Zorn 引理，利用 a)）。
\end{enumerate}

\begin{enumerate}
\def\labelenumi{\arabic{enumi})}
\setcounter{enumi}{11}
\tightlist
\item
  \(\pi\) -准素模 M 是循环子模的直和，当且仅当存在 M 中的拟自由子集 H，使得若 N 是由 H 生成的 M 的纯子模，则 \(M(\pi) = N(\pi)\) 。
\end{enumerate}

\begin{enumerate}
\def\labelenumi{\alph{enumi})}
\setcounter{enumi}{1}
\item
  模 M 是循环子模的直和，当且仅当 M 是子模的某个递增序列 \((P_{k})\) 的并，使得每个 \(P_{k}\) 中非零元素在 M 中的高度是有界的。（要看出该条件是充分的，使用 a) 和习题 11 的 b)）。
\item
  由 \(b)\) 推出，若 \(M\) 是循环子模的直和，则 \(M\) 的每个子模都是循环子模的直和。
\item
  设 M 为一个 \(\pi\) -准素模，没有无限高度元素，具有可数生成系（ \(x_{k}\) ）。证明 M 是循环子模的直和（若
\end{enumerate}

\(P_{k}\) 是由 \(x_{i}\) （给定 \(i \leqslant k\) ）生成的子模，则将 b) 应用于递增序列 \((\mathrm{P}_{k})\) ，并利用习题 5 应用于递减序列 \((\mathrm{P}_{k} \cap \pi^{j}\mathrm{M})\) （对固定的 k））。

\begin{enumerate}
\def\labelenumi{\alph{enumi})}
\setcounter{enumi}{4}
\tightlist
\item
  设 M 为一个 \(\pi\) -准素模，具有可数生成系。则 M 是不可分解子模的直和当且仅当 M 中无限高度元素构成的子模 N 是纯子模（利用 d) 和习题 7，c)）。
\end{enumerate}

\begin{enumerate}
\def\labelenumi{\arabic{enumi})}
\setcounter{enumi}{12}
\tightlist
\item
  设 A 为主理想整环，设 \(\pi\) 为 A 的一个不可约元素，并设 M 为 A 上的 \(\pi\) -准素模。设 I 为满足如下条件的序数 \(\alpha\) （集合论，III，第 225 页，习题14）组成的集：序数之集 \(< \alpha\) 的基数至多等于 M 的基数。由超限归纳法，定义 M 的子模 \(\mathbf{M}^{(\alpha)}\) ——对所有 \(\alpha \in I\) ——如下：取 \(\mathbf{M}^{(0)} = \mathbf{M}\) ；取 \(\mathbf{M}^{(\alpha + 1)} = \pi \mathbf{M}^{(\alpha)}\) ；并且如果 \(\alpha\) 是极限序数，取 \(\mathbf{M}^{(\alpha)}\) 为所有 \(\mathbf{M}^{(\beta)}\) （给定 \(\beta < \alpha\) ）的交集。
\end{enumerate}

\begin{enumerate}
\def\labelenumi{\alph{enumi})}
\item
  证明存在一个最小序数 \(\tau \in I\) 使得 \(M^{(\tau + 1)} = M^{(\tau)}\) （M 的末端序数）。证明 \(M^{(\tau)}\) 在 M 中有补，并且是若干同构于 \(U_{\pi}\) 的子模的直和（习题 3）。称 M 是约化的，如果 \(M^{(\tau)} = 0\) 。
\item
  设 M 为约化 \(\pi\) -模。对所有 \(x \in \mathbf{M}\) ，证明存在一个最大序数 \(\gamma(x) \leqslant \tau\) 使得 \(x \in \mathbf{M}^{(\gamma(x))}\) （特别地， \(\gamma(0) = \tau\) ）。推广习题 8 中的定义，我们称 \(\gamma(x)\) 为 \(x\) 在 M 中的高度。证明若 \(\gamma(x) < \gamma(y)\) ，则 \(\gamma(x + y) = \gamma(x)\) ，且若 \(\gamma(x) = \gamma(y)\) ，则 \(\gamma(x + y) \geqslant \gamma(x)\) 。我们称 \(x\) 关于 M 的子模 N 是适当的，如果对所有 \(y \in \mathbf{N}\) ，高度 \(\gamma(x + y)\) 至多等于 \(\gamma(x)\) ；证明此时 \(\gamma(x + y) = \inf(\gamma(x), \gamma(y))\) 对所有 \(y \in \mathbf{N}\) 成立。
\item
  设 \(\alpha < \tau\) 为一个序数，并设 \(\tilde{\mathbf{M}}^{(\alpha)}\) 为 \(\mathbf{M}^{(\alpha)}\) 的子模，由那些 \(x\) （使得 \(\pi x = \pi y\) 对某个 \(y \in \mathbf{M}^{(\alpha + 1)}\) ）组成。对于那些具有此性质的 \(y \in \mathbf{M}^{(\alpha + 1)}\) ，形如 \(u = x - y\) 的元素构成 模 \(\mathbf{M}(\pi) \cap \mathbf{M}^{(\alpha + 1)}\) 在 \(\mathbf{M}(\pi) \cap \mathbf{M}^{(\alpha)}\) 中的一个类；如果 \(\varphi\) 是从 \(\mathbf{M}(\pi) \cap \mathbf{M}^{(\alpha)}\) 到该模关于 \(\mathbf{M}(\pi) \cap \mathbf{M}^{(\alpha + 1)}\) 的商上的典范同态，则证明映射 \(\psi\) ——它将 \(x\) 映为 \(\varphi(u)\) ——是一个核为 \(\mathbf{M}^{(\alpha + 1)}\) 的同态，并且因此 \(\tilde{\mathbf{M}}^{(\alpha)} / \mathbf{M}^{(\alpha + 1)}\) 同构于 \((\mathbf{M}(\pi) \cap \mathbf{M}^{(\alpha)}) / (\mathbf{M}(\pi) \cap \mathbf{M}^{(\alpha + 1)})\) 。
\item
  设 N 是 M 的一个子模，并设 \(x \in \tilde{\mathbf{M}}^{(\alpha)}\) 是高度为 \(\alpha\) 的元素；证明 x 关于 N 是适当的，当且仅当 \(\mathbf{M}(\pi) \cap \mathbf{M}^{(\alpha)}\) 中对应的元素 u = x - y （参见 c)）关于 N 是适当的。存在这样的元素 x，当且仅当 \(\mathbf{N} \cap \tilde{\mathbf{M}}^{(\alpha)}\) 在同态 \(\psi\) 下的像不同于 \((\mathbf{M}(\pi) \cap \mathbf{M}^{(\alpha)}) / (\mathbf{M}(\pi) \cap \mathbf{M}^{(\alpha + 1)})\) 。
\end{enumerate}

\begin{enumerate}
\def\labelenumi{\arabic{enumi})}
\setcounter{enumi}{13}
\tightlist
\item
  设 M 为约化 \(\pi\) -准素模。在习题 13 的记号中，对于每个序数 \(\alpha < \tau\) ，我们可以把 \((\mathbf{M}(\pi) \cap \mathbf{M}^{(\alpha)}) / (\mathbf{M}(\pi) \cap \mathbf{M}^{(\alpha + 1)})\) 看作域 \(F_{\pi} = A / A\pi\) 上的向量空间；设 \(c(\alpha)\) 表示这个向量空间的维数，并称 \(c(\alpha)\) 为序数 \(\alpha\) 在模 M 中的重数（参见 VII，第 24 页）。
\end{enumerate}

\begin{enumerate}
\def\labelenumi{\alph{enumi})}
\item
  证明：若 M 和 N 是两个 \(\pi\) -准素模，且每个都是循环子模的直和（这意味着 \(\tau \leqslant \omega\) ，即第一个无限序数），则 M 与 N 同构当且仅当它们具有相同的末端序数 \(\tau\) ，并且每个整数 \(n < \tau\) 的重数在 M 中与在 N 中相同。
\item
  给出两个没有无限高度元素的 \(\pi\) -准素模的例子，它们不同构，但其中每个整数 n 具有相同的重数（参见习题 10，f）。c) 设 M 和 N 是两个约化 \(\pi\) -准素模，每个都具有可数无限生成系。证明如果 M 和 N 具有相同的末端序数 \(\tau\) ，并且每个序数 \(\alpha < \tau\) 在 M 和 N 中具有相同的重数，则 M 与 N 同构（« Ulm–Zippin 定理 »）。（给定两个子模 \(S \subset M\) 和 \(T \subset N\) ，一个同构 \(h\) （从 \(S\) 到 \(T\) 上）称为指标同构，如果对所有 \(x \in S\) ， \(x\) 在 \(M\) 中的高度等于 \(h(x)\) 在 \(N\) 中的高度。通过归纳法，并交替交换 \(M\) 和 \(N\) 所扮演的角色，定义一个从 \(M\) 到 \(N\) 上的指标同构，使得问题归结为如下问题：给定两个有限生成的子模 \(S \subset M\) 和 \(T \subset N\) ，一个指标同构 \(h\) （从 \(S\) 到 \(T\) 上），以及一个元素 \(x \in M \cap \complement S\) ，将 \(h\) 延拓为一个从 \(S + Ax\) 到 \(N\) 的一个包含 \(T\) 的子模上的指标同构。证明这可以归结为 \(\pi x \notin S\) 的情形；利用习题 5，选取一个元素 \(w_1\) ，它属于 \((S + Ax) \cap \complement S\) ，具有最大可能的高度 \(\alpha\) （在 \(M\) 中）；然后选取一个元素 \(w_2\) ，它属于 \(S \cap M^{(\alpha)}\) ，使得 \(\pi(w_1 + w_2)\) 在 \(M\) 中的高度是最大可能的；证明 \(w = w_1 + w_2\) 关于 \(S\) 是适当的（习题13，b)）。令 \(z = h(\pi w) \in T\) ；证明，如果存在一个元素 \(t\) （属于 \(N\) ，高度为 \(\alpha\) （在 \(N\) 中），关于 \(T\) 是适当的，且满足 \(\pi t = z\) ），那么 \(\bar{h}\) （ \(h\) 的一个所需形式的扩张）可以通过令 \(\bar{h}(w) = t\) 来得到。剩下的问题是证明元素 \(t\) 的存在性。若
\end{enumerate}

\[
\gamma (z) = \gamma (\pi w) = \alpha + 1 <   \tau ,
\]

即证明存在这样的元素 \(u\) ：它属于 \(\mathbf{N}^{(\alpha)} \cap \complement \mathbf{T}\) ，且 \(z = \pi u\) ，并且 \(t\) 可以取为任意这样的元素。另一方面，如果 \(\gamma(z) = \gamma(\pi w) > \alpha + 1\) ，或者如果 \(\gamma(z) = \gamma(\pi w) = \alpha + 1 = \tau\) ，则利用习题 13， \(d\) ) 来证明元素 \(t\) 的存在性，注意到（在习题 13 的记号下）在 \(\psi\) 之下， \(S \cap \tilde{\mathbf{M}}^{(\alpha)}\) 在 \((\mathbf{M}(\pi) \cap \mathbf{M}^{(\alpha)}) / (\mathbf{M}(\pi) \cap \mathbf{M}^{(\alpha + 1)})\) 中的像是域 \(F_{\pi}\) 上的有限维空间）

* d) 由 c) 得出习题 12 d) 的一个新证明（同时使用第七章第14页的定理1）。 *

\begin{enumerate}
\def\labelenumi{\arabic{enumi})}
\setcounter{enumi}{14}
\tightlist
\item
  \begin{enumerate}
  \def\labelenumii{\alph{enumii})}
  \tightlist
  \item
    设 A 为主理想整环，设 M 为可除 A-模（习题 3），并设 N 为挠 A-模。证明 \(\mathbf{M} \otimes_{\mathbf{A}} \mathbf{N} = 0\) 。
  \end{enumerate}
\end{enumerate}

\begin{enumerate}
\def\labelenumi{\alph{enumi})}
\setcounter{enumi}{1}
\tightlist
\item
  设 M 和 N 为两个没有无限高度元素的 \(\pi\) -准素模；证明 \(M \otimes_{A} N\) 是循环子模的直和（利用 a) 以及习题 10，a)）。
\end{enumerate}

